\chapter{End-to-end workflows}\label{ch:workflows}

Seven workflows exist. Two are the product (W3, W4), one is the research engine (W1) and the rest feed
data into them. Each lists who triggers it, the steps, and what can go wrong.

\section{W1. Reproduce the published numbers}
\textbf{Trigger:} \texttt{make results} (or \texttt{docker run cetp-artifact}). \textbf{Steps:}
(1) \texttt{canonical\_evaluation.py} loads \src{data/corrected/tpch\_dataset\_corrected\_v2.csv}
(2100 rows $=5$ machines $\times$ 21 queries $\times$ 20 repetitions; Q15 excluded). (2) It discards repetition 1 per cell, takes
medians, builds 420 ordered dev$\to$prod pairs. (3) For each of LOMO/LOQO/LOMQO it refits the classifier and the quantile model
inside each fold. (4) It scores six approaches, runs the cluster bootstrap and the leakage audit, then writes
\src{results/canonical\_metrics.json} and \src{results/CANONICAL.md}. \textbf{Check:} \texttt{make verify-determinism}
runs twice and compares. \textbf{Evidence:} the pack author ran this and found byte-identical output \evPACK; I did not \evUNK.
\textbf{Risk:} the Gen-2 script (\texttt{make operator-results}) needs \src{data/plans\_grid}, which is git-ignored.

\section{W2. Capture a machine signature}
\textbf{Trigger:} \texttt{operator\_probe.py --machine-id ID --db DB [--host H --user U]}. \textbf{Steps:} create three \texttt{UNLOGGED}
tables; run seven probes with forced operators; verify the planner used the intended node; compute rates, speed-up, reliability;
write to the registry; drop tables in \texttt{finally}. \textbf{Risks:} a busy machine yields an unreliable record; the registry is a shared
JSON file with no locking; probing production means creating tables there.

\section{W3. Offline prediction from two plan files (\texttt{cetp predict-prod})}
This is the path I ran. \Cref{sec:worked} traces it numerically.
\begin{enumerate}
\item Developer saves \texttt{EXPLAIN (ANALYZE, FORMAT JSON)} on dev, optionally production's plain \texttt{EXPLAIN}.
\item \texttt{cetp predict-prod --dev-plan P --dev D --prod R} loads both machines' rates from the registry
 (refusing v1, invalid or unreliable records), checks probe scales match.
\item The model prices the plan, \texttt{cross\_env.predict} attaches trust and warnings, the CLI prints.
\end{enumerate}
\textbf{Defect:} step 3 omits the machine records, so the cross-ISA refusal is skipped (\cref{sec:cli-bypass}).

\section{W4. Live prediction with plan forcing (\texttt{predict\_prod.py})}
\begin{figure}[ht]
\centering
\begin{tikzpicture}[node distance=5mm and 8mm,
 box/.style={draw=cetpblue,rounded corners=2pt,fill=lightgrey,minimum height=8mm,text width=30mm,align=center,font=\scriptsize},
 prod/.style={draw=cetpred,rounded corners=2pt,fill=cetpred!6,minimum height=8mm,text width=30mm,align=center,font=\scriptsize},
 arr/.style={-{Stealth[length=1.8mm]},thick,cetpgrey}]
\node[prod] (a) {1. Production: \texttt{EXPLAIN} (no ANALYZE), plan only};
\node[box,right=of a] (b) {2. Dev: force that plan with \texttt{pg\_hint\_plan}; default 5 runs, first discarded};
\node[box,right=of b] (c) {3. Verify plan signature matches; take warm median};
\node[box,below=of c] (d) {4. Look up stored calibration for (dev, prod)};
\node[box,left=of d] (e) {5. \texttt{cross\_env.predict}: price nodes};
\node[box,left=of e] (f) {6. Print; log telemetry};
\node[prod,below=of e] (g) {7. Later: \texttt{--actual-ms}, record observation, refit};
\draw[arr] (a)--(b); \draw[arr] (b)--(c); \draw[arr] (c)--(d); \draw[arr] (d)--(e); \draw[arr] (e)--(f); \draw[arr] (f)|-(g);
\draw[arr] (g.east) -| ([xshift=6mm]d.south east) -- (d.south east);
\end{tikzpicture}
\caption{W4, from \src{scripts/collection/predict_prod.py}. Red boxes touch production or the measured truth. Only step 1 contacts production unless \texttt{--actual-ms} is a real measurement.}
\label{fig:w4}
\end{figure}
Unlike W3, this script passes both machine records, so a cross-ISA call with no stored calibration exits with code 4 and prints
``REFUSED''. I did not run it (needs PostgreSQL and \texttt{pg\_hint\_plan}). \textbf{Defect (code reading, not executed):} the refusal handler returns 4 \emph{before} the \texttt{--actual-ms} block runs, so the observation that the message asks the user to supply is never recorded when the store holds no calibration for a cross-ISA pair; the loop cannot bootstrap through this script unless \texttt{--calibration} is set to something other than 1.0. \textbf{Doc mismatch:} the module docstring mentions
\texttt{--calibrate k}, but no such argument exists in the parser \evCODE.

\section{W5. CI gate (Generation 1)}
\texttt{gate/cetp\_gate.py} reads dev features and two Gen-1 signatures, evaluates quantile models and the three-state rule, and exits non-zero on
BLOCK (test: \texttt{test\_block\_exit\_code\_is\_one}). The workflow was changed so BLOCK is reported, not a build failure (commit \texttt{ffce2ec}).

\section{W6. Ingestion and calibration loop}
\texttt{auto\_explain} log $\to$ \texttt{parse\_auto\_explain\_log} (redacted plans + durations) $\to$ \texttt{rank\_queries} (top-3 by total time are
calibration candidates) $\to$ measure those on production $\to$ \texttt{CalibrationStore.record} $\to$ median factor. The learning curve method
shows the factor converging per pair. The loop's only quantitative evidence is a simulation: held-out MAPE falls from 77.72 to 15.01 ($n{=}36$, retrain trigger at 12)
(\evFILE\ \src{results/addition7\_summary.json}), not a deployment.

\section{W7. Container grid collection}
\texttt{container\_grid.py} (\texttt{make grid-setup}, \texttt{grid-collect}) starts \texttt{postgres:15} with CPU in \{1,2,4,8\} and memory in \{1g,3g\},
loads TPC-H, runs the probe and each query for repetitions, and saves plans and timings under \src{data/container\_grid}. It is the only way to
generate Gen-2 evidence and was not run by me.

\section{A worked example: one real execution, step by step}\label{sec:worked}
Task: predict TPC-H q1's runtime on the 1-CPU, 1\,GB container \texttt{c1m1g} from the laptop \texttt{m1-host}. Inputs are the committed
fixture \src{tests/test\_fixtures/q1\_local\_real\_explain.json} and the committed registry. Every number below was produced by running the shipped code
(\evRUN) on the device in this session.

\paragraph{Step 1: decompose.} \texttt{Execution Time} is 3{,}586.94\,ms. The tree has five nodes:
\begin{center}\small
\begin{tabular}{@{}llrrrc@{}}
\toprule \thead{Node} & \thead{Axis} & \thead{Self ms} & \thead{Loops} & \thead{Rows} & \thead{Under Gather}\\ \midrule
Aggregate (Sorted, finalize) & sort & 0.1 & 1 & 4 & no\\
Gather Merge & unmapped & 0.0 & 1 & 12 & no\\
Sort & sort & 0.8 & 3 & 12 & yes\\
Aggregate (Hashed, partial) & hashagg & 3{,}214.5 & 3 & 12 & yes\\
Seq Scan (parallel aware) & seqscan & 7{,}533.2 & 3 & 5{,}927{,}445 & yes\\
\bottomrule
\end{tabular}
\end{center}
Self-times sum to 10{,}748.6\,ms, which is $2.997\times$ the wall time; the \texttt{residual\_frac} is $-1.997$.
The inflation is exactly the averaging effect described in \cref{ch:background}; nothing in the trust logic uses the residual.

\paragraph{Step 2: machine rates (tuples/s) and speed-ups from the registry.}
\begin{center}\small
\begin{tabular}{@{}lrrrr@{}}
\toprule \thead{Machine} & \thead{seqscan} & \thead{hashagg} & \thead{sort} & \thead{parallel speed-up $\pi$}\\ \midrule
m1-host (dev) & 11{,}487{,}849 & 3{,}458{,}930 & 4{,}610{,}457 & 2.6252\\
c1m1g (prod) & 9{,}929{,}772 & 2{,}851{,}544 & 3{,}844{,}003 & 0.7694\\
\bottomrule
\end{tabular}
\end{center}

\paragraph{Step 3: weights.} Parallel buckets are divided by $\pi_d=2.6252$: seqscan $7533.2/2.6252=2869.6$, hashagg
$3214.5/2.6252=1224.5$, sort(par) $0.3$. Total $4094.4$. Weights: seqscan(par) $0.7008$, hashagg(par) $0.2991$, sort(par) $0.0001$,
the serial sort $\approx0$. Unmapped share $0$ ($0.0$\,ms of Gather Merge).

\paragraph{Step 4: ratios.}
$\rho_{\text{seq,par}}=\dfrac{11{,}487{,}849\times2.6252}{9{,}929{,}772\times0.7694}=3.948$;\quad
$\rho_{\text{agg,par}}=\dfrac{3{,}458{,}930\times2.6252}{2{,}851{,}544\times0.7694}=4.139$.

\paragraph{Step 5: scaling.} $s=0.7008\times3.948+0.2991\times4.139+\ldots=4.0049$ (program output). Predicted runtime $=3586.9\times4.0049=14{,}365$\,ms.

\paragraph{Step 6: sanity check against stored measurements.} The fixture is one particular run. The committed evaluation row for the same query and
machine used a different dev run (2{,}246.7\,ms) and records predicted scaling 4.008 (vs 4.005 here: weights differ slightly between runs),
predicted 9{,}005.6\,ms against an actual 7{,}874.0\,ms, i.e.\ $+14.4\%$ error versus $71.5\%$ for the naive baseline \evFILE\ (\src{results/laptop\_to\_prod\_pairs.csv}).
This is the ``good'' case.

\paragraph{Step 7: the same call towards an x86 runner.} For \texttt{gha-x86} the program gives $s=1.8717$ (seqscan(par) ratio 1.8745, hashagg(par) 1.8652):
the model reports the target about $1.87\times$ slower. The measured value for the committed row is $1.239$: predicted 4{,}204.8\,ms vs actual 2{,}782.8\,ms ($+51.1\%$), naive error
19.3\%. Over the three queries q1/q18/q19 on this runner the operator model errs by 51\%, 57\% and 80\%, naive by 19\%, 6\% and 2.5\% \evFILE. This is the cross-ISA failure in miniature.

\paragraph{Step 8: what the CLI printed.} \texttt{scaling factor: 1.8717}, \texttt{predicted prod: 6,713.6 ms} (for the fixture's 3{,}586.9\,ms dev run),
\texttt{explained: 100\%}, \texttt{trust: MEDIUM}, and four warnings, one of which cites stale numbers (\cref{sec:cli-bypass}). Exit status 0.

\begin{notebox}
The same example by hand for Gen 1 is \cref{ch:algo}: c7i$\to$c5a, $\sigma=0.5917$, 3{,}587\,ms $\to$ 2{,}122\,ms.
\end{notebox}

\section{The CLI bypasses the cross-architecture refusal}\label{sec:cli-bypass}
\textbf{Finding (\evRUN).} The library function \texttt{cross\_env.predict} refuses an uncalibrated cross-ISA prediction \emph{only if it is given the two machine
records}. The \texttt{cetp predict-prod} command (\src{cetp/cli.py}:146--190) loads the records into \texttt{dev\_rec}/\texttt{prod\_rec} but does not pass them. Result:
\begin{lstlisting}[language={},mathescape=false,caption={Reproduction (read-only; no repository file was modified).}]
$ python3 -m cetp.cli predict-prod --dev m1-host --prod gha-x86 \
      --dev-plan tests/test_fixtures/q1_local_real_explain.json
... scaling factor: 1.8717   predicted prod: 6,713.6 ms   trust: MEDIUM   (exit 0)

>>> cross_env.predict(None, plan, 2246.7, dr, pr, dev_rec=drec, prod_rec=prec)
CrossArchitectureUncalibrated: refusing an uncalibrated prediction from arm64 to x86_64 ...
\end{lstlisting}
\textbf{Severity:} high for the product claim (``refuses what it measured to be worse than doing nothing''), low for the research claims. \textbf{Why tests missed it:} no test invokes
\texttt{cmd\_predict\_prod}. \textbf{Fix (not applied):} add \texttt{dev\_rec=dev\_rec, prod\_rec=prod\_rec} to the call and a CLI test.
